% =========================================================================
% SEÇÃO 6: CONCLUSÕES
% =========================================================================
\section{Conclusões}\label{sec:conclusoes}

Os resultados obtidos demonstram o cumprimento integral dos objetivos do plano de trabalho, articulando teoria matemática, engenharia de software em C++23 e validação experimental:

\begin{enumerate}[leftmargin=1.5em, itemsep=0.2em, topsep=0.2em, parsep=0em]
    \item \textbf{Fundamentação Teórica:} Demonstração formal dos teoremas centrais da teoria de fluxos: Fluxo Máximo e Corte Mínimo, Integralidade de Ford-Fulkerson, Decomposição de Fluxo, Teorema de Menger (arcos e vértices), Teorema de König e as Quatro Dualidades em grafos bipartidos. No domínio de custo mínimo: condição de otimalidade por ciclos negativos de Klein, equivalência dos custos reduzidos sob potenciais nodais, complementaridade de folgas KKT e Total Unimodularidade da matriz de incidência;

    \item \textbf{Arquitetura em C++23:} Desenvolvimento de classes organizadas em duas famílias polimórficas com representação de adjacência contígua e reversão em $\mathcal{O}(1)$:
    \begin{itemize}[leftmargin=1.5em, itemsep=0.1em, topsep=0.1em, parsep=0em]
        \item \textbf{Família \lstinline{FlowNetwork} (Fluxo Máximo):} Centralizada na classe base abstrata \lstinline{FlowNetwork} e implementada pelos algoritmos \lstinline{FordFulkerson}, \lstinline{EdmondsKarp}, \lstinline{Dinic}, \lstinline{PushRelabel} (FIFO) e \lstinline{PushRelabelImproved} (Gap Heuristic), documentados nos Apêndices~\ref{ap:flownetwork} a~\ref{ap:push_relabel_improved};
        \item \textbf{Família \lstinline{CostNetwork} (Custo Mínimo):} Centralizada na classe base abstrata \lstinline{CostNetwork} e implementada pelos algoritmos \lstinline{CycleCanceling}, \lstinline{SuccessiveShortest} (SPFA), \lstinline{SuccessiveShortestDijkstra} (Potenciais) e \lstinline{NetworkSimplex}, documentados nos Apêndices~\ref{ap:costnetwork} a~\ref{ap:network_simplex};
    \end{itemize}

    \item \textbf{Reduções Combinatórias e Aplicações:} Resolução de problemas clássicos (\textit{Maximum Bipartite Matching}, \textit{Minimum Vertex Cover}, \textit{Maximum Independent Set}, caminhos disjuntos e desconexão de redes) e implementação de \textbf{Segmentação Interativa de Imagens} baseada em corte mínimo via corte de grafos (\textit{Graph Cuts});

    \item \textbf{Avaliação Experimental:} Protocolo experimental com instâncias da literatura (DIMACS Washington, Genrmf e Netgen), utilizando isolamento estrito de tempo de CPU e validação cruzada com 100\% de concordância entre os algoritmos implementados.
\end{enumerate}

Como trabalhos futuros, destacam-se a avaliação empírica sobre instâncias esparsas massivas, a análise de diferentes regras de pivoteamento no \textit{Network Simplex} (como \textit{Candidate List Pricing}) e a preparação de artigo científico com a síntese dos resultados.
